% Source: https://artofproblemsolving.com/community/c3h2027730_geometric_or_algebric
% Stable identity: aops-topic-2027730-q01
During second $k$ ($k=1,\dots,40$) the bug first walks $1$ mm outward, reaching distance $k$ mm from the centre; then the tick carries it along an arc of radius $k$ and angle $6^\circ=\dfrac{\pi}{30}$, of length $\dfrac{k\pi}{30}$.
\begin{asy}
usepackage("amsmath");
size(5cm);
defaultpen(fontsize(10pt));
pair O=(0,0);
real th=6*pi/180;
pair p=O;
for(int k=1;k<=5;++k){
  pair q=k*dir(90-6*(k-1));
  draw(p--q,blue+linewidth(0.9));
  pair q2=k*dir(90-6*k);
  draw(arc(O,k,90-6*(k-1),90-6*k),red+linewidth(0.9));
  p=q2;
}
for(int k=0;k<=5;++k) draw(O--5.6*dir(90-6*k),gray+0.3);
dot(O); label("$O$",O,SW);
label("walk $1$",(0.5*dir(90))+(-0.05,0),W,blue);
label("tick: arc $\dfrac{k\pi}{30}$",4.6*dir(70),NE,red);
label("\mbox{}",(min(currentpicture,user=true)+max(currentpicture,user=true))/2,black); // keeps the text after the figure black when LaTeX typesets the labels inline
\end{asy}

The total distance is
\[
40+\dfrac{\pi}{30}(1+2+\dots+40)=40+\dfrac{820\pi}{30}=40+\dfrac{82\pi}{3}=\dfrac{120+82\pi}{3}.
\]
Here $\gcd(120,82,3)=1$, so $a=120$, $b=82$, $c=3$ and
\[a+b+c=\boxed{205}.\]
